% !TeX root = ../sections/02_exercises.tex
\subsection{1_1}
\begin{exercise}[情報数学 A 中間試験 1]
	周期 \(2\pi\) をもち，区間 \([-\pi,\pi)\) 上で
	\[
	\frac{1}{x},\quad \sqrt{|x|},\quad x,\quad x^2,\quad \sin 3x
	\]
	で与えられる五つの周期関数について，次の条件に当てはまるものを全て挙げよ．
	なお，該当する関数が存在しない場合は「なし」と答えよ．
	
	\medskip
	
	※解答の際に，連続性の証明（確認）は不要で，単に分類結果のみを記せばよい．
	
	\begin{enumerate}
		\item 連続な関数
		\item 区分的に連続な関数
		\item 滑らかな関数
		\item 区分的に滑らかな関数
	\end{enumerate}
\end{exercise}
\begin{background}
	この問題では，区間 \([-\pi,\pi)\) 上で与えられた関数を，
	周期 \(2\pi\) の周期関数として考える。
	
	したがって，単に区間の内部での性質を見るだけでは不十分である。
	周期 \(2\pi\) で延長するということは，点 \(-\pi\) と点 \(\pi\) が
	つながるという意味であるから，端点でのつながり方も確認する必要がある。
	
	この問題で判定すべき性質は，次の四つである。
	
	まず，連続な関数とは，周期関数として全ての点で連続な関数である。
	特に，区間内部で連続であることに加えて，
	\[
	\lim_{x\to \pi-0}f(x)=f(-\pi)
	\]
	となっているかを確認する必要がある。
	
	次に，区分的に連続な関数とは，有限個の点を除いて連続であり，
	不連続点においても左右極限が有限に存在する関数である。
	したがって，ジャンプ型の不連続は許されるが，
	\[
	\frac{1}{x}
	\]
	のように，ある点で無限大に発散する関数は区分的に連続とはいえない。
	
	次に，滑らかな関数とは，周期関数として何回でも微分可能な関数である。
	そのため，関数の値だけでなく，導関数，高階導関数まで端点でつながっている必要がある。
	
	最後に，区分的に滑らかな関数とは，有限個の点を除いて滑らかであり，
	各区間では十分よく微分できる関数である。
	ただし，授業での定義によっては，導関数の片側極限の有限性まで要求する場合がある。
	そのため，\(\sqrt{|x|}\) のように折れ曲がり，あるいは導関数の発散をもつ関数は
	定義に照らして慎重に判断する必要がある。
	
	この問題では，次の観点が重要である。
	
	\[
	\begin{array}{c|c}
		\text{確認する場所} & \text{見るべき性質} \\ \hline
		\text{区間内部} & \text{定義されているか，連続か，微分可能か} \\
		\text{端点} & \text{周期拡張で値や導関数がつながるか} \\
		\text{特異点} & \text{無限大に発散しないか} \\
		\text{折れ曲がり点} & \text{滑らかさに影響するか}
	\end{array}
	\]
\end{background}
\begin{strategy}
	この問題を解くときは，次の順番で各関数を調べる。
	
	\begin{enumerate}
		\item まず，区間 \([-\pi,\pi)\) の内部で定義されているかを確認する。
		\item 次に，内部で連続かどうかを見る。
		\item 次に，周期拡張したときに端点で値がつながるかを見る。
		\item さらに，導関数が存在するか，また端点で導関数がつながるかを見る。
		\item 最後に，連続，区分的連続，滑らか，区分的滑らかのどれに入るかを分類する。
	\end{enumerate}
	
	まず，
	\[
	\frac{1}{x}
	\]
	は \(x=0\) で定義されず，さらに \(x\to 0\) で無限大に発散する。
	したがって，単なるジャンプ不連続ではない。
	そのため，連続でも区分的に連続でもないと判断する。
	
	次に，
	\[
	\sqrt{|x|}
	\]
	は \(x=0\) でも値が定義でき，関数としては連続である。
	また，端点でも
	\[
	\sqrt{|-\pi|}=\sqrt{\pi}
	\]
	であり，周期的な値のつながりも問題ない。
	しかし，\(x=0\) で尖った形をしており，導関数が通常の意味ではよく振る舞わない。
	したがって，連続性と滑らかさを分けて判断する必要がある。
	
	次に，
	\[
	x
	\]
	は区間内部では滑らかである。
	しかし周期拡張すると，\(\pi\) の左側では値が \(\pi\) に近づく一方，
	\(-\pi\) では値が \(-\pi\) である。
	つまり端点でジャンプが生じる。
	したがって，連続ではないが，ジャンプ不連続は有限個なので，
	区分的に連続，区分的に滑らかであると判断する。
	
	次に，
	\[
	x^2
	\]
	は区間内部で滑らかであり，端点でも
	\[
	(-\pi)^2=\pi^2
	\]
	となるので，周期拡張しても値はつながる。
	したがって連続である。
	しかし導関数は
	\[
	2x
	\]
	であり，端点で
	\[
	\lim_{x\to \pi-0}2x=2\pi,
	\qquad
	\lim_{x\to -\pi+0}2x=-2\pi
	\]
	となる。
	よって導関数は周期的につながらない。
	したがって，滑らかではないが，区分的に滑らかである。
	
	最後に，
	\[
	\sin 3x
	\]
	は三角関数であり，もともと滑らかである。
	さらに \(2\pi\) 周期で見ても値や導関数が自然につながる。
	したがって，滑らかな関数である。
	
	この問題で最も大切なのは，
	「区間内部でよい関数」かどうかだけでなく，
	周期関数として端点がつながっているかを必ず確認することである。
\end{strategy}